\documentclass[11pt]{article}
\usepackage[a4paper,margin=26mm]{geometry}
\usepackage{amsmath,amssymb,mathtools}
\usepackage{unicode-math}
\setmainfont{Latin Modern Roman}
\setmathfont{Latin Modern Math}
\usepackage{xurl,hyperref}
\hypersetup{hidelinks,pdftitle={Problem statement}}
\newcommand{\sref}[2]{\hyperlink{source:#1:#2}{[S#2]}}
\newcommand{\cataloguescope}[1]{#1}
\begin{document}
\section{Planar half-density percolation has no unique infinite cluster}\label{statement}

\subsection{Definitions and mathematical statement}

\paragraph{Shared notation.}$\mathbb N=\{1,2,\ldots\}$, and $\mathbb Z,\mathbb Q,\mathbb R,\mathbb C$ denote the integers, rationals, reals and complex numbers. Any other conventions are specified locally.


Let $G=(V,E)$ be a connected locally finite planar graph. In independent site percolation with parameter $p$, each vertex is open with probability $p$, independently of all others. An open cluster is a connected component of the subgraph induced by open vertices. Write $N_\infty$ for the number of infinite open clusters and $\mathbb P_p$ for this product probability law.
\paragraph{Statement recorded in the supplied source.}
At $p=1/2$, does
\[\mathbb P_{1/2}(N_\infty\ge1)=1\quad\Longrightarrow\quad\mathbb P_{1/2}(N_\infty=\infty)=1\]
hold for every such planar graph? This is the site-percolation question, with no transitivity or minimum-degree hypothesis. \sref{statement}{2}
\subsection{Short English statement}
Does almost-sure site percolation at probability one half on a planar graph force infinitely many infinite open clusters?
\subsection{Sources}
\begin{description}
\item[\hypertarget{source:statement:1}{S1}] ULAM.AI and the UnsolvedMath Contributors, \emph{UnsolvedMath: A Curated Collection of Open Mathematics Problems}, catalogue entry. CC BY 4.0. \url{https://huggingface.co/datasets/ulamai/UnsolvedMath}.
\item[\hypertarget{source:statement:2}{S2}] I. Benjamini and O. Schramm, \emph{Percolation beyond $\mathbb Z^d$, many questions and a few answers}, Electronic Communications in Probability 1 (1996), pp. 71--82, Conjecture 8, p. 80, with the conventions of Section 2. \url{https://emis.de/ft/41630}.
\end{description}
\label{end:statement}
\end{document}
